% 393_Rauschgrenze_Lernende_Systeme_De_ch.tex
% Johann Pascher, 13. Okt. 2026
% Rauschgrenze und lernende Systeme: Fragen und Antworten

\providecommand{\BM}{\textbf{[B]}}
\providecommand{\KM}{\textbf{[K]}}
\providecommand{\SM}{\textbf{[S]}}
\providecommand{\XM}{\textbf{[X]}}
\providecommand{\QM}{\textbf{[Q]}}

\chapter*{Rauschgrenze und lernende Systeme}
\addcontentsline{toc}{chapter}{Rauschgrenze und lernende Systeme}

\begin{abstract}
Der Korpus kennt eine Rauschgrenze: Ab einer bestimmten Messtoleranz sind harmonische Zusammenhänge
nicht mehr von Zufall zu unterscheiden (Dok.~343, Satz G$'$). Dieses Dokument geht in Frage und
Antwort der Frage nach, ob neuronale Netze und Sprachmodelle an dieselbe Grenze stoßen, warum sie
trotzdem fast immer lesbare Sätze bilden, warum Sprache Fehler verzeiht und Zahlen nicht, und welche
Rolle Wahrscheinlichkeit und die Gewichte eines Netzes dabei spielen. Ein Nulltest mit vollständiger
symbolischer Suche über rund 23 Millionen Ausdrücke zeigt, dass die dimensionslosen Messgrößen der
Grundform allein keine harmonische Struktur über Zufallsniveau tragen, und dass selbst eine echte
Regel wie die Balmer-Serie aus wenigen Beispielen nur mit der passenden, engen Vorgabe sichtbar wird.
Der gemeinsame Kern aller Antworten: Diskrete Struktur ist fehlertolerant, die Rauschgrenze sitzt
am Übergang vom Diskreten zum Kontinuum, und Struktur, die die Daten nicht hergeben, muss vorgegeben
und dann geprüft werden.
\end{abstract}

% ============================================================
\section*{Statusmarkierungen}
% ============================================================
\begin{center}
\begin{tabular}{@{}lp{11cm}@{}}
\textbf{[B]} & Algebraisch oder informationstheoretisch bewiesen. \\[2pt]
\textbf{[K]} & Numerisch bestätigt (Prüfskript). \\[2pt]
\textbf{[S]} & Setzung, offen oder Einordnung ohne eigene Rechnung. \\[2pt]
\textbf{[X]} & Durch die Rechnung ausgeschlossen. \\[2pt]
\textbf{[Q]} & Aus der Literatur übernommen, mit Quelle. \\
\end{tabular}
\end{center}

\noindent Das Dokument ist als Folge von Fragen und Antworten aufgebaut. Die Fragen sind in der
Reihenfolge belassen, in der sie sich auseinander ergeben; jede Antwort steht für sich. Abschnitt~14
fasst den gemeinsamen Kern in einem Bild zusammen.

% ============================================================
\section{Die Rauschgrenze im Korpus}
% ============================================================

\textbf{Frage.} Was ist mit der Rauschgrenze gemeint?

\textbf{Antwort.} Ein gemessenes Verhältnis ist nur bis auf eine Toleranz $\varepsilon$ bekannt. Je
größer der zugelassene Nenner, desto dichter liegen die Brüche; ab der Farey-Höhe
\[
  Q_c=\frac{\pi}{\sqrt{6\,\varepsilon}}
\]
liegen im Toleranzfenster so viele Brüche, dass ein echtes Verhältnis nicht mehr von einem zufälligen
zu trennen ist (Dok.~343, Satz G$'$). Bei $\varepsilon=1{,}05\,\%$, der Größe der fraktalen Korrektur,
ist $Q_c\approx12{,}5$; bei $\varepsilon=100\,\xi\approx1{,}33\,\%$ ist $Q_c\approx11{,}1$ (Dok.~359,
Satz~B) \KM. Harmonische Zusammenhänge mit kleinen Nennern bleiben sichtbar, feinere gehen im Rauschen
unter. Die Grenze ist eine Eigenschaft der Daten, nicht des Verfahrens: Was im Signal nicht enthalten
ist, holt keine Analyse zurück, weder eine Fourier-Analyse noch eine Autokorrelation \BM.

% ============================================================
\section{Stoßen neuronale Netze und Sprachmodelle an dieselbe Grenze?}
% ============================================================

\textbf{Frage.} Müssten neuronale Netze und Sprachmodelle bei der Mustererkennung nicht ebenfalls an
diese Grenze stoßen?

\textbf{Antwort.} Ja. Ein Netz ist ein Verfahren wie jedes andere; es kann die Grenze erreichen, aber
nicht überschreiten. Jenseits der Grenze erkennt es nicht mehr, sondern passt sich an das Rauschen an
und hält zufällige Nähe für Struktur; in der Lerntheorie heißt das Überanpassung. Hinzu kommt eine
bekannte Eigenschaft solcher Netze: Sie lernen grobe, niederfrequente Strukturen zuerst und feine
spät \QM\,\cite{Rahaman2019}. Feine Strukturen sind aber gerade die, die dem Rauschen am nächsten
liegen.

Sprachmodelle sind in einer ungünstigeren Lage. Sie lernen aus Text, nicht aus Messsignalen, und haben
keine eigene Messauflösung. Bei Zahlenverhältnissen neigen sie zu Falschtreffern: Bei genügend vielen
Kombinationen findet sich fast immer eine \enquote{harmonische} Erklärung (Look-elsewhere-Effekt).
Deshalb werden im Korpus Prüfskripte als Nulltests angelegt: Liefert eine Zufallsauswahl dasselbe
Ergebnis, wird es nicht gewertet. Das gilt für jeden Vorschlag eines Sprachmodells.

Besser als die Rohdaten werden Netze nur durch eingebautes Vorwissen. Dok.~359 Satz~D zeigt das an
einem Beispiel: Ob das Zulässigkeitskriterium \enquote{Primfaktoren nur aus $\{2,3,5,7,11,13\}$}
lernbar ist, entscheidet die Hypothesenklasse. Ohne Strukturannahme verlangt die Lernschranke rund
84\,000 markierte Beispiele, mit der Strukturannahme \enquote{Primfaktoren aus einer festen Menge}
höchstens 407; aus positiven Beispielen allein wird die richtige Primzahlmenge dann schon nach 30
Beispielen sicher erkannt \KM\,\cite{ShalevShwartz2014}.

% ============================================================
\section{Warum bilden Sprachmodelle trotzdem lesbare Sätze?}
% ============================================================

\textbf{Frage.} Wenn das so ist, warum erzeugen Sprachmodelle kaum je unlesbare Sätze?

\textbf{Antwort.} Weil Grammatik und Satzbau das sind, was in den Trainingsdaten am dichtesten
vorkommt. Es sind Billionen Wörter, dieselben Regeln wiederholen sich in Millionen verschiedener
Sätze, es gibt kein Messrauschen, und die Architektur ist auf Abfolgen von Wörtern zugeschnitten. Das
ist genau der günstige Fall aus Dok.~359 Satz~D: passende Vorgabe und reichlich Beispiele.

Schwach werden die Modelle dort, wo das Muster selten oder exakt ist: genaue Zahlen, Rechenwege,
seltene Fakten, Literaturangaben. Dort erzeugen sie Plausibles, und der Satz bleibt dabei flüssig.
Lesbarkeit ist das dichte Muster, Richtigkeit oft das dünne. Ein gut geschriebener Satz eines
Sprachmodells sagt deshalb nichts über die Richtigkeit seines Inhalts.

% ============================================================
\section{Müssten harmonische Muster nicht so dicht sein wie Grammatik?}
% ============================================================

\textbf{Frage.} Wenn die FFGFT zutrifft, durchziehen harmonische Verhältnisse die ganze Natur. Müssten
sie dann nicht ebenso dicht vorkommen wie Grammatik?

\textbf{Antwort.} Dicht im Vorkommen ja, dicht in der Vielfalt nein. Für das Lernen einer Regel zählt
die Zahl \emph{verschiedener} Beispiele, nicht die Zahl der Wiederholungen. Das Elektron kommt
unvorstellbar oft vor, liefert aber jedes Mal dasselbe Massenverhältnis. Die Grundstruktur der FFGFT
steckt in wenigen Zahlen: einigen Massenverhältnissen, Kopplungen und Mischungswinkeln, zusammen ein
paar Dutzend. Aus so wenigen Werten kann kein Verfahren eine Regel von Zufall trennen.

Wo harmonische Muster tatsächlich in großer Vielfalt vorkommen, wurden sie auch gefunden: in
Atomspektren mit Tausenden Linien nach einem ganzzahligen Bauprinzip, in Kristallgittern, in
Bahnresonanzen und Musikinstrumenten. Balmer fand die Regel der Wasserstofflinien 1885 aus vier
Linien von Hand \QM\,\cite{Balmer1885}; symbolische Regression hat bekannte Formeln der Physik aus
Daten wiedergefunden \QM\,\cite{Udrescu2020}.

Ein zweiter Grund kommt aus dem Korpus selbst. Die gemessenen Verhältnisse sind nach FFGFT keine
exakten Brüche: Die Grundform ist exakt, der Messwert trägt den Einbettungsfehler (Dok.~333, Dok.~358
Satz~A), und die PDG-Werte enthalten modellabhängige Anteile (Dok.~351). Ein Lernverfahren, das nur
Messwerte sieht, lernt die verschobene Form, nicht die Regel darunter. In der Sprache gibt es diese
Schicht nicht.

% ============================================================
\section{Der Nulltest}
\label{sec:nulltest}
% ============================================================

\textbf{Frage.} Lässt sich prüfen, wie viel Struktur die Messwerte selbst tragen und wie viel die
Vorgabe beiträgt?

\textbf{Antwort.} Ja, mit einer symbolischen Regression in ihrer exakten Form, der vollständigen
Suche über eine Ausdrucks-Grammatik. Durchsucht werden alle Ausdrücke
\[
  x=c\cdot d^{1/k}\cdot\xi^{m/n}\cdot\pi^{s},
\]
mit teilerfremdem Bruch $c=p/q$, $p,q\le200$, Wurzeln $d^{1/k}$ mit $d\le30$ und $k\in\{2,3\}$,
$\xi=1/7500$ mit $n\in\{1,2,3,6\}$ und $|m/n|\le4{,}5$, sowie $s\in\{-1,0,1\}$. Die Primfaktoren von
$p$, $q$ und $d$ stammen aus einer gewählten Primzahlmenge. Mit $\{2,3,5,7,11,13\}$ sind das
23\,084\,325 Ausdrücke; getrennt wird nur über die Brüche $p/q$ gesucht (3109 Ausdrücke). Jeder
Ausdruck erhält eine Komplexität in Bit, die mit Zähler, Nenner, Wurzel und Exponent wächst. Der
$p$-Wert ist der Anteil zufälliger Zahlen (log-gleichverteilt, eine Dekade um den Messwert), die die
Grammatik bei gleicher oder geringerer Komplexität ebenso genau trifft. Treffer, die genauer sind als
die Messunsicherheit, zählen nur bis zu dieser.

\textbf{Mit Vorgabe.} Die Formeln der Grundform aus Dok.~387 sind fest; verglichen wird mit allen
gleich komplexen Ausdrücken derselben Grammatik:

\begin{center}
\small
\begin{tabular}{@{}lllr@{}}
\toprule
\textbf{Größe} & \textbf{Formel} & \textbf{Abweichung} & \textbf{$p$} \\
\midrule
$v/E_P$ & $\xi^4/(5\pi)$ & $-0{,}23\,\%$ & $0{,}02$ \\
$M_W/M_Z$ & $\sqrt{7/9}$ & $+0{,}06\,\%$ & $0{,}09$ \\
$\lambda_{\mathrm{CKM}}$ & $\xi^{1/6}$ & $+0{,}45\,\%$ & $0{,}11$ \\
$m_h/M_Z$ & $11/8$ & $+0{,}15\,\%$ & $0{,}42$ \\
$m_\mu/m_e$, $m_\tau/m_\mu$, $\Delta m^2$-Verhältnis, $\sin^2\theta_{13}$, $m_t/M_Z$ &  & $0{,}1$--$1{,}5\,\%$ & $1{,}0$ \\
\bottomrule
\end{tabular}
\end{center}

\noindent Zusammengefasst nach Fisher ergibt das $p=0{,}41$ \KM. Bei $0{,}5$ bis $1\,\%$ Abweichung
trifft die Grammatik jede Zahl; die Leptonformeln haben dort numerisch keine Beweiskraft. Am
auffälligsten ist $v/E_P$. Für $\alpha^{-1}=3700/27$ (Dok.~338) gilt: Die Abweichung von $7{,}6$\,ppm
erreicht ein Bruch mit Nenner bis 27 auch bei einer Zufallszahl in etwa der Hälfte der Fälle
($p\approx0{,}47$) \KM. Das gilt allerdings nur, wenn man $3700/27$ als frei gesuchten Bruch
behandelt. Im Korpus entsteht er als Kürzung von $\tfrac{74}{75}\cdot7500/54$ aus $1/\xi$, der Regel
$74/75$ und der ganzen Zahl $54$; nach diesen Bestandteilen gerechnet liegt die
Zufallswahrscheinlichkeit bei rund $10^{-3}$ (Dok.~394, 395) \KM.

\textbf{Ohne Vorgabe.} Die blinde Suche über dieselben neun Größen bleibt für jede Primzahlmenge auf
Zufallsniveau: für $\{2,3,5,7,11,13\}$, für die Kontrollmengen $\{2,3,5,7,17,19\}$ und
$\{2,3,5,13,17,23\}$, für $\{2,3,5,7\}$ und ebenso für reine Zufallszahlen \KM.

\textbf{Positivkontrolle.} Die Balmer-Linien des Wasserstoffs haben die exakten Verhältnisse
$27/20$, $28/25$ und $200/189$; die Luftwellenlängen treffen sie auf $10^{-5}$ \QM\,\cite{NIST}. Mit
der engen Grammatik der Brüche wird diese Regel schon aus drei Verhältnissen gefunden (zusammen
$p=0{,}002$), mit der vollen Grammatik nicht ($p=0{,}47$) \KM. Die 23 Millionen Ausdrücke liegen so
dicht, dass auch Zufallszahlen ebenso gut getroffen werden.

\textbf{Was daraus folgt.} Selbst eine wahre Regel ist aus wenigen Beispielen nur mit der richtigen,
engen Vorgabe sichtbar. Die numerische Nähe trägt die FFGFT nicht; das muss die Ableitungskette tun:
Geometrie, Galois-Struktur, Formel. Das entspricht der Lesart von Dok.~384 und Dok.~387, nach der die
Verhältnisse der Grundform die Struktur prüfen. Widerlegt ist durch den Test nichts; er misst nur,
was die Zahlen allein hergeben.

\textbf{Grenzen des Tests.} Das Bit-Maß der Komplexität ist eine Wahl; ein anderes Maß verschiebt die
einzelnen $p$-Werte \SM. Es sind nur neun Größen. Die Formeln sind in Kenntnis der Messwerte
entstanden; der Vergleich mit gleich komplexen Alternativen ist deshalb der faire Maßstab, nicht der
rohe Treffer. Die Leptonreste von $0{,}5$ bis $1\,\%$ gelten im Korpus als echte Struktur
($K_{\mathrm{frak}}$, Dok.~352); der Test kann sie nur als Fehltreffer werten.

% ============================================================
\section{Sprache ändert sich laufend, und trotzdem funktioniert es}
% ============================================================

\textbf{Frage.} Sprache ändert sich dauernd, es entstehen neue Wortbildungen, Grammatik variiert in
Dialekten. Warum funktioniert es trotzdem?

\textbf{Antwort.} Gerade die Veränderlichkeit bestätigt das Bild. Sprache arbeitet mit großer
Toleranz und liefert zugleich sehr viele Beispiele.

Zur Toleranz: Ein Satz bleibt verständlich, wenn ein Wort falsch geschrieben, ein Dialektwort
eingesetzt oder die Satzstellung verdreht ist. Gedruckte Sprache ist stark redundant
\QM\,\cite{Shannon1951}. Die Genauigkeit, die Sprache braucht, ist grob; in der Sprache des Korpus:
großes $\varepsilon$, damit kleines $Q_c$. Sprache lebt in dem Bereich, in dem Strukturen noch
trennbar sind.

Zu den Beispielen: Neue Wortbildungen und Dialekte sind keine Einzelfälle, sondern selbst Muster mit
vielen Belegen. Ein Dialekt hat Millionen Sprecher und unzählige Sätze; neue Wörter folgen alten
Bauregeln, \enquote{googeln} wird konjugiert wie \enquote{segeln}. Ein Lernverfahren lernt deshalb
keine einzelne starre Grammatik, sondern eine Verteilung von Varianten, und jede Variante ist dicht
belegt. Unter den Dialekten liegt eine gemeinsame tiefere Struktur, wie bei Balmer: viele Linien,
eine Regel. Wo es nicht funktioniert, passt es ebenfalls: Bei einem kleinen Dialekt mit wenig Text,
einem seltenen Fachwort oder einem ganz neuen Ausdruck werden Sprachmodelle unsicher.

Die Musik zeigt beide Seiten. Die harmonische Grammatik aus Oktave, Quinte und Terz, also aus den
Primzahlen 2, 3 und 5, ist universell erkennbar und hat viele \enquote{Dialekte}: reine Stimmung,
gleichstufig, mitteltönig, Volksmusik-Stimmungen. Sie funktioniert, weil das Ohr Abweichungen in der
Größe eines Kommas toleriert. Das ist die Aussage von Dok.~342 und Dok.~343: Bis $k\approx3$ sind die
Galois-Strukturen bei dieser Toleranz trennbar, darüber nicht mehr \BM. Die Harmonik der kleinen
Primzahlen ist in der Natur so dicht und robust wie Grammatik. Die feinen Verhältnisse der
Teilchenmassen, bei denen es auf Zehntelprozent ankommt, liegen jenseits dieser Grenze; dort muss die
Struktur aus der Ableitung kommen.

% ============================================================
\section{Warum verzeiht Sprache Fehler und Zahlen nicht?}
% ============================================================

\textbf{Frage.} Warum lässt Sprache diese Toleranz zu, Mathematik und Zahlen aber nicht?

\textbf{Antwort.} Wegen der Redundanz. Nachrichtentechnisch gesprochen ist Sprache ein
fehlerkorrigierender Code, eine Zahl eine uncodierte Übertragung.

Die meisten zufälligen Buchstabenfolgen sind keine Wörter. Wird ein Buchstabe falsch, entsteht meist
ein ungültiges Wort, und der Empfänger korrigiert zum nächsten gültigen; Satzbau und Zusammenhang
prüfen zusätzlich. Das entspricht einem Code mit großem Hamming-Abstand: Fehler werden erkannt und
korrigiert. Viele Formen führen außerdem zur selben Bedeutung, und am Ziel zählt nur die Bedeutung.

Bei einer Zahl ist der Abstand null. Jede Ziffernfolge ist eine gültige Zahl; aus $137{,}036$ wird
bei einem Ziffernfehler $137{,}037$, keine fehlerhafte, sondern eine andere, genauso gültige Zahl.
Nichts zeigt den Fehler an, nichts korrigiert ihn. Mathematik verkettet außerdem: Ein Fehler geht in
jeden folgenden Schritt ein und wird oft verstärkt, während sich Fehler in einem Satz nicht
aufsummieren.

Beim Zählen ist Mathematik dagegen tolerant. Ganze Zahlen, Gruppenordnungen und Wicklungszahlen
verhalten sich wie Wörter: Zwischen 3 und 4 gibt es keine gültige Gruppenordnung, und $3{,}01$ wird
sofort als ungültig erkannt. Im Korpus steht das als \enquote{Zustand gegen Prozess}: Zählen ist
exakt, ein Bit ist eine Wicklungszahl (Dok.~343, Dok.~247). Die algebraisch bewiesenen Aussagen der
FFGFT, also $\mathbb Z_3$, $GF(27)$, Orbits und Fixpunkte, sind deshalb so fehlertolerant wie
Sprache. Die Intoleranz entsteht an einer einzigen Stelle: beim Übergang vom Diskreten zum
Kontinuum, wenn ein ganzzahliges Verhältnis mit einem Messwert auf der reellen Achse verglichen wird.
Dort ist jeder Nachbarwert gültig, und es gibt keinen Abstand, der das Rauschen abfängt.

% ============================================================
\section{Schöpfen KI-Modelle aus einem fertigen Rahmenwerk?}
% ============================================================

\textbf{Frage.} KI-Modelle haben letztlich ein fertiges Rahmenwerk, aus dem sie schöpfen; es muss
nicht alles neu nach Mustern durchsucht werden. Wie verträgt sich das mit der Rauschgrenze?

\textbf{Antwort.} Nach dem Training ist ein Modell tatsächlich ein fertiges Rahmenwerk; beim Antworten
wird nichts mehr durchsucht, das Gelernte wird angewendet. Die Gewichte sind die eingefrorene
Vorgabe, im Sinn von Dok.~359 Satz~D die Hypothesenklasse.

Dieses Rahmenwerk ist aber selbst aus Daten gelernt. Es ist eine verdichtete Form der Trainingstexte
und enthält nur, was dort dicht vorkam. Wo die Texte dünn waren, hat es Lücken, und das Modell füllt
sie mit Plausiblem; daher kommen erfundene Literaturangaben. Es ist außerdem statistisch und nicht
exakt, für Sprache gebaut und für exakte Zahlen unzuverlässig. Moderne Systeme rufen deshalb für den
exakten Teil ein Werkzeug auf: Das statistische Rahmenwerk entscheidet, was zu rechnen ist, das
Werkzeug rechnet. Bei der Arbeit am Korpus ist das die Arbeitsteilung zwischen Formulierung und
Prüfskript.

Auch die FFGFT ist ein fertiges Rahmenwerk: $\xi$, $T^4/\mathbb Z_3$, die Galois-Struktur. Es ist
aber nicht aus Daten gelernt, sondern gesetzt und abgeleitet. Das gelernte Rahmenwerk ist breit und
tolerant, aber ohne Garantie für Exaktheit; das gesetzte ist exakt und diskret, muss sich aber am
Messwert bewähren, genau an der Stelle, wo Diskretes auf Kontinuum trifft. Ein Lernverfahren kann
die FFGFT-Struktur nach dem Nulltest nicht selbst finden; es kann sie aber als Rahmen eingebaut
bekommen, wie es die Architekturskizze von Dok.~359 vorschlägt \SM.

% ============================================================
\section{Woher kommen die Regeln eines KI-Assistenten?}
% ============================================================

\textbf{Frage.} Ein KI-Assistent arbeitet mit vielen vorgegebenen Regeln. Sind diese gelernt oder
direkt eingegeben?

\textbf{Antwort.} Beides, in drei Schichten.

Die erste Schicht ist gelernt und macht den größten Teil aus. Zuerst entsteht beim Training an sehr
vielen Texten das statistische Rahmenwerk. Danach wird in einem zweiten Training das Verhalten
geformt, mit Bewertungen durch Menschen und mit schriftlich formulierten Grundsätzen wie Ehrlichkeit,
Schadensvermeidung und dem Eingeständnis von Unsicherheit. Auch diese Grundsätze stecken danach in den
Gewichten und werden nicht bei jeder Antwort nachgelesen.

Die zweite Schicht wird bei jedem Gespräch neu mitgegeben: Anweisungen des Produkts, etwa wie
Werkzeuge benutzt werden und was nur nach Rückfrage geschehen darf, sowie gespeicherte Vorgaben des
Nutzers. Diese Schicht ist am leichtesten zu ändern; sie gilt sofort, ohne neues Training. Die dritte
Schicht ist der laufende Zusammenhang: alles, was im Gespräch steht und was gelesen wurde.

Die Schichten können sich widersprechen. Bei der Arbeit an diesem Korpus verlangte ein automatischer
Prüfschritt nach jeder Antwort, Änderungen selbst in das öffentliche Repositorium zu übertragen,
während die gespeicherte Vorgabe des Autors genau das untersagt; die ausdrückliche Vorgabe wurde
befolgt. Zum Thema dieses Dokuments passt die Aufteilung: Die erste Schicht ist das gelernte,
tolerante Rahmenwerk, die zweite und dritte sind gesetzte Regeln, die Prüfskripte das exakte
Werkzeug, das beides kontrolliert.

% ============================================================
\section{Wie entsteht nach einer Fehlermeldung der nächste Schritt?}
% ============================================================

\textbf{Frage.} Zwischen einer Fehlermeldung eines Skripts und der Entscheidung, was als Nächstes am
sinnvollsten ist, müssen doch viele Schichten liegen.

\textbf{Antwort.} Ja, aber anders, als man vermuten würde: Es gibt keinen eigenen Entscheidungsbaustein,
sondern Schichten im Netz und eine einfache Schleife außen herum.

Die Schleife: Der Assistent schreibt einen Werkzeugaufruf. Das umgebende Programm führt ihn aus und
hängt das Ergebnis, auch eine Fehlermeldung, an den Gesprächstext an. Dann wird das Ganze neu gelesen
und der nächste Schritt geschrieben. Ein Plan existiert nur als Text in diesem Zusammenhang.

Die Schichten: Das Netz selbst hat Dutzende bis über hundert hintereinanderliegende Schichten, von
denen jede den ganzen bisherigen Text weiterverarbeitet. Aus Fehlermeldung, Ziel, gelesenen Dateien
und bisherigen Versuchen entsteht so der nächste Schritt. Dass er sinnvoll ist, ist trainiert: Im
zweiten Training wurden solche Abläufe an vielen Aufgaben geübt, und Abläufe, die zum Ziel führten,
wurden verstärkt.

Zwei Beispiele aus dem Neusatz der Buchreihe: Eine Meldung \enquote{Improper alphabetic constant}
beim Setzen eines Kapiteltitels mit $\varphi$ ist ein dicht belegtes Muster; die Ursache, Mathematik
in den PDF-Lesezeichen, wird sofort erkannt. Ein Band, der ohne Fehlermeldung nur 18 statt über 300
Seiten ergab, verlangte dagegen einen Schluss: Die Eingabe endete zu früh, vermutlich durch ein
\verb|\end{document}| in einem Kapitel, was eine Suche bestätigte. Wo die Muster dünn sind, wird
probiert, und es braucht mehrere Anläufe. Deshalb gilt für den Assistenten dasselbe wie für die
FFGFT: ausprobieren, exakt prüfen, mit Prüfskripten nachmessen, statt dem eigenen Eindruck zu trauen.

% ============================================================
\section{Zufall und Wahrscheinlichkeit}
% ============================================================

\textbf{Frage.} Spielt bei Zufall und Wahrscheinlichkeiten die Rauschgrenze nicht ebenfalls eine
Rolle?

\textbf{Antwort.} Ja, sogar ganz direkt: Die Wahrscheinlichkeitsrechnung ist die Mathematik der
Rauschgrenze. Dabei sind zwei Arten von Wahrscheinlichkeit zu unterscheiden.

Wahrscheinlichkeit als Nichtwissen: Ein Messwert schwankt. Mittelt man $N$ unabhängige Messungen,
sinkt das Rauschen wie $1/\sqrt N$; die hundertfache Zahl von Messungen senkt es auf ein Zehntel \BM.
Deshalb ist $m_\mu/m_e$ auf $10^{-8}$ bekannt, obwohl jede einzelne Messung gröber ist. Das gilt nur
für zufällige Fehler; systematische Anteile wie die Modellannahmen der PDG-Werte (Dok.~351) mitteln
sich nicht weg. Auch der $p$-Wert des Nulltests gehört hierher: Er ist die Rauschgrenze in
Wahrscheinlichkeit ausgedrückt.

In Sprachmodellen ist Wahrscheinlichkeit die Ausgabe selbst. Bei jedem Wort wird eine Verteilung über
alle möglichen nächsten Wörter berechnet und eines ausgewählt. Wo die Trainingsdaten dicht waren, ist
die Verteilung scharf und die Antwort stabil; wo sie dünn waren, ist sie flach, und die Auswahl
entscheidet beinahe der Zufall. Eine flache Verteilung ist die Rauschgrenze des Modells, und von
außen ist sie nicht zu sehen, weil der Satz trotzdem flüssig klingt.

Wahrscheinlichkeit als Struktur: In der Quantenmechanik sind Wahrscheinlichkeiten exakte Zahlen. Nach
Dok.~367 sind die Koide-Phase $\theta$ und die ikosaedrische Übergangswahrscheinlichkeit $p_0$
derselbe Quotient $2/3^2=2/9$, zwei Lesarten einer Born-Regel-Struktur. Eine solche Wahrscheinlichkeit
ist so diskret wie eine Gruppenordnung. Rauschen kommt erst bei der Messung hinzu: Eine Häufigkeit
aus $N$ Ereignissen streut um $2/9$ mit etwa $\sqrt{p(1-p)/N}$; um $2/9$ auf drei Standardabweichungen
vom Nachbarbruch $1/5$ zu trennen, braucht es rund 3200 Ereignisse \KM. In diesem Sinn führt
Dok.~230 auch die Bell-Korrelationen: real, aber geometrisch strukturiert. Wahrscheinlichkeit kommt
damit zweimal vor: als Gegenstand der Theorie, exakt, und als Werkzeug, um zu beurteilen, ob das
Exakte im Rauschen noch erkennbar ist.

% ============================================================
\section{Fließt die Rauschgrenze in die Gewichte ein?}
% ============================================================

\textbf{Frage.} Wie ist das mit den Gewichten eines Netzes: Fließen dort die Rauschgrenzen ein?

\textbf{Antwort.} Ja, auf zwei Wegen, und die Gewichte sind selbst ein Beispiel für das Prinzip.

Beim Lernen werden die Gewichte in vielen kleinen Schritten angepasst, jeder nach einem kleinen Stapel
von Beispielen und deshalb verrauscht. Was oft vorkommt, schiebt die Gewichte immer wieder in dieselbe
Richtung und setzt sich durch; das Rauschen mittelt sich mit $1/\sqrt N$ weg. Was selten vorkommt,
geht im Rauschen der Schritte unter. Dazu kommen gezielt eingebaute Bremsen: Gewichte werden klein
gehalten, und das Training wird rechtzeitig beendet. Damit wird bewusst eine Rauschgrenze gesetzt:
Feinstruktur, die nicht oft genug belegt ist, soll nicht gelernt werden, weil sie meist Rauschen ist.

Beim Speichern ist jedes einzelne Gewicht erstaunlich ungenau. Mit 16 Bit sind es etwa $2{,}4$ bis
$3{,}3$ Dezimalstellen \KM, und oft wird nachträglich auf 8 oder 4 Bit gekürzt, ohne dass das Modell
merklich schlechter wird. Das geht nur, weil das Wissen nicht in einzelnen Gewichten steckt, sondern
verteilt in sehr vielen zugleich. Das Netz ist damit selbst ein redundanter, fehlertoleranter Code,
wie die Sprache. Gute Trainingslösungen haben genau diese Eigenschaft: Eine kleine Verschiebung der
Gewichte ändert das Ergebnis kaum \QM\,\cite{Hochreiter1997}. Die Kehrseite ist dieselbe wie bei der
Sprache: Exakte Rechnungen, die an jeder Stelle volle Genauigkeit verlangen, gelingen im Inneren
schlecht.

% ============================================================
\section{Funktioniert ohne gesetzte Rahmenbedingungen überhaupt etwas?}
% ============================================================

\textbf{Frage.} Heißt das, dass ohne gesetzte Rahmenbedingungen ohnehin nichts funktioniert?

\textbf{Antwort.} Im Kern ja; für das Lernen aus Daten ist das bewiesen. Nach dem No-Free-Lunch-Theorem
der Lerntheorie ist über alle denkbaren Aufgaben gemittelt kein Lernverfahren besser als ein anderes,
auch nicht besser als Raten \QM\,\cite{Wolpert1996}. Ein Verfahren ist nur dann besser, wenn seine
eingebauten Annahmen zur tatsächlichen Welt passen. Jedes Verfahren hat solche Annahmen, auch wenn sie
nicht ausgesprochen werden: Die Fourier-Analyse setzt Periodizität voraus, ein Bildnetz, dass
benachbarte Bildpunkte zusammengehören, ein Sprachmodell, dass Text eine Abfolge mit Struktur ist.
Dok.~359 Satz~D zeigt den Unterschied in Zahlen: ohne Rahmen rund 84\,000 Beispiele, mit dem
passenden Rahmen höchstens 407 und aus positiven Beispielen allein 30 \KM. Ohne Vorgabe ist jedes
Muster so gut wie jedes andere.

Zwei Einschränkungen gehören dazu. Erstens kann ein Rahmen schwach sein, wenn Daten reichlich
vorhanden sind: Je mehr Daten, desto weniger muss vorgegeben werden. Sprachmodelle haben einen recht
allgemeinen Rahmen und gleichen ihn durch sehr große Textmengen aus. Je dünner die Daten, desto stärker
und genauer muss der Rahmen sein; bei wenigen Dutzend Naturkonstanten muss er praktisch die ganze
Struktur tragen. Das ist die Lage der FFGFT.

Zweitens funktioniert ein Rahmen nicht schon deshalb, weil er gesetzt ist. Auch ein falscher oder zu
weiter Rahmen passt sich an und findet dann falsche Muster. Der Balmer-Test aus
Abschnitt~\ref{sec:nulltest} zeigt beide Seiten: Der enge Rahmen der Brüche findet die wahre Regel;
der weite Rahmen mit 23 Millionen Ausdrücken passt auf alles und unterscheidet dadurch nichts mehr
\KM. Ein Rahmen muss deshalb nicht nur gesetzt, sondern auch geprüft werden; der Nulltest misst, ob
er mehr leistet als Zufall.

Für die Physik heißt das: Jede Theorie hat Setzungen. Das Standardmodell hat seine Eichgruppen und
rund zwei Dutzend gemessene Parameter, die FFGFT hat $T^4/\mathbb Z_3$, $\xi$ und die ganzzahligen
Strukturwahlen. Die Frage ist nicht, ob es Setzungen gibt, sondern wie wenige es sind, wie viel aus
ihnen folgt und ob das Gefolgerte einer Prüfung gegen Zufall standhält: Setzung, Ableitung,
Vergleich, Nulltest.

% ============================================================
\section{Der gemeinsame Kern in einem Bild}
% ============================================================

Alle Antworten laufen auf dieselbe Struktur hinaus. Auf der diskreten Seite stehen Zählen, Wörter,
Gruppenordnungen, Wicklungszahlen und Born-Gewichte; dort gibt es Abstände zwischen den gültigen
Werten, und Rauschen, das kleiner ist als dieser Abstand, wird abgefangen. Auf der kontinuierlichen
Seite stehen Messwerte und die Gewichte eines Netzes; dort ist jeder Nachbarwert gültig.

\begin{center}
\begin{tikzpicture}[
  box/.style={draw, rounded corners, align=center, minimum width=5.2cm, minimum height=1.6cm, font=\small},
  arr/.style={-{Stealth[length=3mm]}, thick}]
\node[box] (D) at (0,0) {\textbf{Diskret}\\Zählen, Wörter, Gruppenordnungen,\\Wicklungszahlen, Born-Gewicht $2/9$\\\emph{Abstand fängt Rauschen ab}};
\node[box] (K) at (8,0) {\textbf{Kontinuum}\\Messwerte, Gewichte eines Netzes\\[2pt]\emph{jeder Nachbarwert gültig}};
\draw[arr] (D.north east) to[bend left=18] node[above, font=\small, align=center]{Messung: Struktur wird abgebildet\\\textbf{hier sitzt die Rauschgrenze}} (K.north west);
\draw[arr] (K.south west) to[bend left=18] node[below, font=\small, align=center]{Ausgabe eines Netzes: Wort wird gewählt\\Rauschen unter dem Abstand wird abgeschnitten} (D.south east);
\end{tikzpicture}
\end{center}

\noindent Ein Netz rechnet im Inneren kontinuierlich und gibt diskrete Wörter aus; beim Schritt vom
Kontinuum zum Diskreten wird das Rauschen abgeschnitten, solange es kleiner ist als der Abstand
zwischen zwei Wörtern. Bei der Messung läuft es umgekehrt: Eine exakte, diskrete Struktur wird auf das
Kontinuum abgebildet, und dort fängt kein Abstand das Rauschen ab. Daraus folgen die vier Kernsätze
dieses Dokuments:

\begin{enumerate}
\item Diskrete Struktur ist fehlertolerant; die Rauschgrenze sitzt am Übergang vom Diskreten zum
  Kontinuum \BM.
\item Struktur, die in den Daten dicht und vielfältig belegt ist, wird gelernt; Struktur aus wenigen
  exakten Werten wird aus Daten nicht gefunden, auch von einem Netz nicht \KM.
\item Solche Struktur muss vorgegeben und dann geprüft werden: Setzung, Ableitung, Vergleich mit dem
  Messwert, Nulltest. Diesen Weg geht die FFGFT \SM.
\item Ohne Rahmen kein Lernen, aber ein Rahmen ist nur so gut wie seine Prüfung \QM\,\cite{Wolpert1996}, \KM.
\end{enumerate}

\vspace{1em}
\noindent\textit{Prüfskripte:} \texttt{2/python/Dok393\_Skripte/} ---
\texttt{pruef\_393a\_nulltest\_grammatik.py} (8/8, Ergebnisprotokoll
\texttt{pruef\_393a\_ergebnis.txt}), \texttt{pruef\_393b\_zahlen.py} (11/11); zusammen
\texttt{run\_all\_393.sh}, 19/19.

\section*{Werkzeug und Methode}
\addcontentsline{toc}{section}{Werkzeug und Methode}

Claude (Anthropic) wurde als Werkzeug eingesetzt für: Ausformulieren der Antworten nach den Fragen des
Autors, Schreiben und Ausführen der Prüfskripte und Ausarbeiten der \LaTeX-Quellen. Die Abschnitte
über die Arbeitsweise von KI-Assistenten beschreiben die allgemeine Funktionsweise solcher Systeme; sie
beruhen nicht auf einer Einsicht in die eigenen Gewichte. Jede Zahl im Text ist durch eine Prüfung in
den genannten Skripten abgesichert; Literaturangaben sind als \QM\ markiert.

\begin{thebibliography}{99}
\bibitem{Balmer1885} J.~J.~Balmer, \textit{Notiz über die Spectrallinien des Wasserstoffs}, Annalen der Physik \textbf{261} (1885) 80--87.
\bibitem{Shannon1951} C.~E.~Shannon, \textit{Prediction and entropy of printed English}, Bell System Technical Journal \textbf{30} (1951) 50--64.
\bibitem{Hochreiter1997} S.~Hochreiter, J.~Schmidhuber, \textit{Flat minima}, Neural Computation \textbf{9} (1997) 1--42.
\bibitem{Rahaman2019} N.~Rahaman u.\,a., \textit{On the spectral bias of neural networks}, Proceedings of the 36th International Conference on Machine Learning, PMLR \textbf{97} (2019) 5301--5310.
\bibitem{Udrescu2020} S.-M.~Udrescu, M.~Tegmark, \textit{AI Feynman: A physics-inspired method for symbolic regression}, Science Advances \textbf{6} (2020) eaay2631.
\bibitem{Wolpert1996} D.~H.~Wolpert, \textit{The lack of a priori distinctions between learning algorithms}, Neural Computation \textbf{8} (1996) 1341--1390.
\bibitem{ShalevShwartz2014} S.~Shalev-Shwartz, S.~Ben-David, \textit{Understanding Machine Learning: From Theory to Algorithms}, Cambridge University Press (2014).
\bibitem{NIST} A.~Kramida, Yu.~Ralchenko, J.~Reader, NIST ASD Team, \textit{NIST Atomic Spectra Database}, National Institute of Standards and Technology, Gaithersburg (Wasserstofflinien H$\alpha$--H$\delta$, Luftwellenlängen).
\end{thebibliography}
